%% ma: L12-B03-0002
%% lop: 12
%% chuong: 1
%% bai: 3
%% dang: 12.3.1
%% loai: TN4
%% muc: TH
%% dap_an: "D"
%% nguon: "(NBV)-ĐỀ SỐ 1-MỖI NGÀY 1 ĐỀ THI 2025.pdf (D01, MỖI NGÀY 1 ĐỀ - Đề số 1), Phần I Câu 2"
%% kien_thuc: "Bài 3 (tiệm cận ngang của đồ thị hàm số)"
%% ghi_chu: "Trùng ý tưởng với câu cùng dạng, giáo viên đã duyệt giữ (09/10/2026). Hình dựng lại từ hàm y=(1-x)/(x+1) (tiệm cận đứng x=-1, tiệm cận ngang y=-1) theo hình gốc"
%% trang_thai: chinh-thuc
\begin{ex}
Cho hàm số $y=f(x)$ có đồ thị như hình vẽ. Đường thẳng nào sau đây là tiệm cận ngang của đồ thị hàm số đã cho?
\begin{center}
\begin{tikzpicture}[x=1cm,y=1cm]
  \begin{scope}
    \clip (-3.2,-3.4) rectangle (3.2,3.2);
    \draw[thick,domain=-3.1:-1.2,samples=60,smooth] plot (\x,{(1-\x)/(\x+1)});
    \draw[thick,domain=-.67:3.1,samples=80,smooth] plot (\x,{(1-\x)/(\x+1)});
  \end{scope}
  \draw[phu] (-1,-3.4)--(-1,3.2);
  \draw[phu] (-3.2,-1)--(3.2,-1);
  \draw[truc] (-3.2,0)--(3.4,0) node[below]{$x$};
  \draw[truc] (0,-3.5)--(0,3.4) node[left]{$y$};
  \draw (-1,.06)--(-1,-.06) node[above left]{\footnotesize$-1$};
  \draw (.06,-1)--(-.06,-1) node[left,xshift=-2pt,yshift=-3pt]{\footnotesize$-1$};
  \path (0,0) node[below left]{\footnotesize$O$};
\end{tikzpicture}
\end{center}
\choice{$x=1$}{$x=-1$}{$y=1$}{\True $y=-1$}
\loigiai{Khi $x\to\pm\infty$ đồ thị tiến sát đường thẳng nằm ngang cắt trục tung tại điểm có tung độ $-1$. Vậy tiệm cận ngang là $y=-1$ (đường $x=-1$ là tiệm cận đứng).}
\end{ex}
